\section{A five-primary sieve at thirteen and the complete fixed rational locus}
\label{sec-fixed-rational-locus}

The ordinary arguments QL, HT, LH, ZS, DS, UD, SU, ZD and IF, followed
by the exact sieve below, prove
\[
 C(\mathbb Q)=\{(1,8),(1,-8),(-1,8),(-1,-8),\infty_+,\infty_-\}.
\]
Here $C$ is the smooth projective completion of
$W^2=z^6-27z^4+99z^2-9$, with $W/z^3=\pm1$ at its two infinities.
This is a theorem for one fixed curve. Its complete height and local
analytic dependency chain has not been formalized in Lean.

\paragraph{SM1: transport through the global rational group.}
Let $P_1=(-2,1)$ on $E_1$ and $P_2=(6,9)$ on $E_2$.
QL2 proves that $E_i(\mathbb Q)$ is torsion-free of rank one and
$d_i=[E_i(\mathbb Q):\mathbb ZP_i]$ is finite and prime to five.
QL1 gives nonzero logarithms of valuation one. For any
$Q\in E_i(\mathbb Q)$ put $m_i(Q)=\log_{E_i}(Q)/\log_{E_i}(P_i)$.
If the good reduction at thirteen has order fifteen, then
\[
 m_i(Q)\in\mathbb Z_5,\qquad
 [3]\operatorname{red}_{13}(Q)
 =[m_i(Q)\bmod5]\,[3]\operatorname{red}_{13}(P_i).
 \tag{SM1}
\]
To prove this, the finite quotient of order $d_i$ gives an integer
$a$ with $[d_i]Q=[a]P_i$. Apply the logarithm to get $m_i(Q)=a/d_i$.
Reduce the integer group relation at thirteen and multiply by three.
Both resulting points are killed by five, so multiplication by $d_i$
is invertible on them; its inverse modulo five proves (SM1).
The residue of $a/d_i$ is exactly that of the five-adic coefficient.
This does not reduce a value of $\mathbb Q_5$ modulo thirteen.
It needs neither a global generator assertion $d_i=1$ nor saturation
at thirteen.

\paragraph{SM2: the complete obstruction over $\mathbb F_{13}$.}
There is no $X\in C(\mathbb F_{13})$ satisfying
\[
 [3]R_1(X)=O,\qquad [3]R_2(X)\in\{(11,2),(11,11)\}.
 \tag{SM2}
\]
Both reduced elliptic groups have order fifteen, and
$[3]\operatorname{red}_{13}(P_2)=(11,2)$.
Here is an exhaustive finite proof. The two discriminants are
respectively $3$ and $9$ modulo thirteen. Their complete affine point
lists, in addition to $O$, are
\[
\begin{split}
 E_1:&\ (0,\pm2),(1,\pm3),(3,\pm2),(6,\pm6),
       (10,\pm2),(11,\pm1),(12,\pm5),\\
 E_2:&\ (3,\pm2),(5,\pm6),(6,\pm4),(8,\pm5),
       (9,\pm1),(11,\pm2),(12,\pm2).
\end{split}
\]
Testing the thirteen possible abscissas and ordinates proves exhaustion
and hence order fifteen in each case. The reduced sextic is
$F=z^6+12z^4+8z^2+4$, and the polynomial identity
\[
 (10+8z^2)F+(6z+3z^3)F'=1\quad\text{in }\mathbb F_{13}[z]
\]
proves squarefreeness even over the algebraic closure. The monic
even-degree model is smooth at its two infinities too.
The complete projective quotient table is
\begin{center}
\begin{tabular}{c|c|c}
 $X$ on $C$ & $[3]R_1(X)$ & $[3]R_2(X)$\\\hline
 $(0,2)$ & $(3,11)$ & $O$\\
 $(0,11)$ & $(3,2)$ & $O$\\
 $(1,5)$ & $(3,11)$ & $(11,2)$\\
 $(1,8)$ & $(3,2)$ & $(11,11)$\\
 $(3,3)$ & $(11,1)$ & $(11,11)$\\
 $(3,10)$ & $(11,12)$ & $(11,2)$\\
 $(6,3)$ & $O$ & $O$\\
 $(6,10)$ & $O$ & $O$\\
 $(7,3)$ & $O$ & $O$\\
 $(7,10)$ & $O$ & $O$\\
 $(10,3)$ & $(11,1)$ & $(11,2)$\\
 $(10,10)$ & $(11,12)$ & $(11,11)$\\
 $(12,5)$ & $(3,11)$ & $(11,11)$\\
 $(12,8)$ & $(3,2)$ & $(11,2)$\\
 $\infty_+$ & $O$ & $(8,8)$\\
 $\infty_-$ & $O$ & $(8,5)$
\end{tabular}
\end{center}
The fourteen affine points exhaust the same thirteen-by-thirteen
enumeration. At $z=0$ the projective quotient values are
$R_2=O$, $R_1=(-9/4,W/8)$; at infinity of sign $\epsilon$ they are
$R_1=O$, $R_2=(33/4,-9\epsilon/8)$. Applying chord and tangent
addition gives the displayed triples and $[3]P_2$.
The six rows with $[3]R_1=O$ establish (SM2).

The reduction table applies to every rational point, including points
outside the integral affine chart. If $z$ is thirteen-adically integral,
then $W$ is integral. If $z$ reduces to zero, $R_2$ has negative
$x$-valuation and reduces to $O$, while $R_1$ is regular. If $z$ has
negative valuation, put $q=1/z$ and $v=Wq^3$: then $q$ reduces to zero
and $v$ to $\pm1$, with the stated infinity quotient values. In the
remaining cases all displayed affine denominators are units. Thus
reduction of the rational quotient images agrees with the table.

Three independent finite implementations reproduce the full table,
using chord and tangent arithmetic, division polynomials, and cubic
line intersection. The primary canonical certificate also contains
the smoothness identity, both entire elliptic point lists and both
forbidden signs. Finite computation here proves only these finite inputs.

\paragraph{SM3: exclude the four extra local zeros.}
The complete ordinary disk certificates give exactly ten distinct
$\mathbb Q_5$ zeros of $f$, counted with total multiplicity twelve:
\begin{center}
\begin{tabular}{c|c|c}
 disk orbit & distinct zeros & multiplicities\\\hline
 $z\equiv\pm1$ & $8$ & all simple\\
 $z\equiv\pm2$ & $0$ & none\\
 $z\equiv0$ & $0$ & none\\
 infinity & $2$ & both double
\end{tabular}
\end{center}
DS covers all twelve residue disks. Four of the unit-disk zeros are
the known rational points $(\pm1,\pm8)$. The other four are the
orbit of the extra UD zero on $z=1+5s$, $W(1)=8$, with $s\equiv4$.
On that whole disk UD gives
\[
 \ell_1/5\equiv4+4s,\qquad \ell_2/5\equiv3+s\pmod5,
\]
while QL gives $\log_{E_1}(P_1)/5\equiv4$ and
$\log_{E_2}(P_2)/5\equiv2$. The extra zero therefore has
$(m_1(R_1),m_2(R_2))\equiv(0,1)\pmod5$. Under DS, either both
quotient points change sign or only the second does. All four extra
zeros thus have coefficients $(0,\pm1)$ modulo five.
If one were rational, (SM1) would force its reduced quotients to
satisfy (SM2), a contradiction. None of the four is rational.

Every finite nonzero rational point satisfies $f=0$ by LH and QL7.
The zero fiber has no rational point because $W^2=-9$; the infinity
points were treated directly in IF2. The list is therefore exhaustive.
Conversely the sextic takes value $64$ at $z=\pm1$ and its leading
coefficient is one, so all six displayed rational points exist.
This completes the fixed-curve ordinary proof. The subsequent exact
inverse tests the second square at the same source and the real
positive domain of the original problem. No varying-curve or
varying-residual uniform height theorem follows from this classification.
